\documentclass[11pt]{article}
\usepackage[margin=1in]{geometry}
\usepackage{amsmath,amssymb,amsfonts,mathtools}
\usepackage{array,booktabs,enumitem,graphicx,xcolor,hyperref}
\usepackage[T1]{fontenc}
\usepackage[utf8]{inputenc}
\title{Classical Dynamics of Point Particles in 2+1 Gravity}
\author{Source-backed arXiv sample}
\date{}
\begin{document}
\maketitle
\section{Figure Captions}

\begin{figure}[htbp]
\centering
\fbox{\rule{0pt}{1.15in}\rule{0.78\linewidth}{0pt}}
\caption{(a) Minkowski space with excised region, and (b) its tail representation in the \textbackslash{}CS solution. The trajectory of a test particle is also indicated.}
\label{fig:source-1}
\end{figure}


\begin{figure}[htbp]
\centering
\fbox{\rule{0pt}{1.15in}\rule{0.78\linewidth}{0pt}}
\caption{The particle exchange operation \$\textbackslash{}sigma\_\{12}
\label{fig:source-2}
\end{figure}


\begin{figure}[htbp]
\centering
\fbox{\rule{0pt}{1.15in}\rule{0.78\linewidth}{0pt}}
\caption{Sequence of exchanges for the Yang-Baxter equation.}
\label{fig:source-3}
\end{figure}


Figures~\ref{fig:source-1} through~\ref{fig:source-3} preserve source-backed captions while replacing external graphics with deterministic placeholders.
\end{document}
